\subsection{同型分量}

定义 9. --- 设 G 是拓扑群，\(\pi\) 是 G 的连续不可约表示。设 \(\varrho\) 是 G 在分离局部凸空间 E 中的连续表示。称 \(\pi\)-同型分量（对应 \(\varrho\)）为 \(\varrho\) 的所有与 \(\pi\) 同构的子表示的表示空间之和在 E 中的闭包。该子空间记为 \(\mathrm{M}_{\pi}(\varrho)\)。

空间 \(\mathrm{M}_{\pi}(\varrho)\) 是 E 的闭子空间，并定义 \(\varrho\) 的一个子表示。该空间只取决于 \(\pi\) 在 \(\widehat{G}\) 中的类。

命题 8. — 设 G 是拓扑群，\(\varrho\) 是 G 在复希尔伯特空间 E 中的酉表示。

\begin{enumerate}
\def\labelenumi{\alph{enumi})}
\item
  设 \(\pi_1\) 和 \(\pi_2\) 是 G 的非同构不可约酉表示。子空间 \(\mathrm{M}_{\pi_1}(\varrho)\) 和 \(\mathrm{M}_{\pi_2}(\varrho)\) 正交；
\item
  设 \(\varrho'\) 是 G 的酉表示，且 \(u\) 是从 \(\varrho\) 到 \(\varrho'\) 的 G-态射。对于 G 的每个不可约酉表示 \(\pi\)，有 \(u(\mathrm{M}_{\pi}(\varrho)) \subset \mathrm{M}_{\pi}(\varrho')\)。
\end{enumerate}

设 \(E_{1}\) 和 \(E_{2}\) 是 E 的子空间，分别定义与 \(\pi_{1}\) 和 \(\pi_{2}\) 同构的子表示。E 上像为 \(E_{2}\) 的正交投影通过限制到子空间定义了 \(\mathrm{Hom}_{\mathrm{G}}(\pi_{1},\pi_{2})\) 的一个元素；该元素为零（V, p. 387 的推论 2），这证明了 \(E_{2}\) 与 \(E_{1}\) 正交。断言 \(a\) 得证。

对于断言 \(b\)，注意到根据定义，\(\mathrm{M}_{\pi}(\varrho)\) 是 E 中由如下元素生成的空间的闭包：元素 \(x \in \mathrm{E}\) 属于闭子空间 \(\mathrm{F} \subset \mathrm{E}\)，该子空间在 \(\varrho\) 下稳定，且其子表示 \(\varrho_{\mathrm{F}}\) 是 \(\varrho\) 的子表示并与 \(\pi\) 同构。因此只需证明在这种情形下 \(u(x) \in \mathrm{M}_{\pi}(\varrho')\)。令 \(\mathrm{H} = u(\mathrm{F})\)；它是 \(\varrho'\) 的表示空间的闭子空间（V, p. 388 的推论 5），在 \(\varrho'\) 下稳定，而由限制到子空间得到的 u 是从 F 到 H 的满射 G-态射。如果 H 非零，则由 V, p. 387 的推论 2，该 G-态射是同构。因此 G 在 H 上的表示与 \(\pi\) 同构，从而 \(u(x)\) 属于 \(\mathrm{M}_{\pi}(\varrho')\)。

对于任意向量空间 H 及 H 的子空间族 \((\mathrm{H}_{i})_{i\in\mathrm{I}}\)，称这些空间 \((\mathrm{H}_{i})_{i\in\mathrm{I}}\) 为直和，如果族 \((\mathrm{H}_{i})_{i\in\mathrm{I}}\) 满足 A, II, p. 18 的命题 11 的等价条件。

命题 9. --- 设 G 是拓扑群，\(\varrho\) 是 G 在复希尔伯特空间 E 中的酉表示。设 \(\pi\) 是 G 在复希尔伯特空间 \(\mathrm{E}_{\pi}\) 中的不可约酉表示。记 \(v\) 为从 \(\mathrm{Hom}_{\mathrm{G}}(\pi, \varrho) \otimes \mathrm{E}_{\pi}\) 到 E 的线性映射，满足 \(v(u \otimes x) = u(x)\) 对所有 \((u, x) \in \mathrm{Hom}_{\mathrm{G}}(\pi, \varrho) \times \mathrm{E}_{\pi}\) 成立。

线性映射 \(v\) 是单射，其像是 \(\varrho\) 的所有与 \(\pi\) 同构的子表示的表示空间之和。特别地，\(v\) 的像在 \(M_{\pi}(\varrho)\) 中稠密。

V, p. 388 的推论 5 表明，v 的像是 \(\varrho\) 的所有与 \(\pi\) 同构的子表示的表示空间之和。

证明 \(v\) 是单射。设 \((u_i)_{i \in I}\) 是向量空间 Hom \(_G(\pi, \varrho)\) 的一组基。对于 \(i \in I\)，记 F \(_i\) 为 \(u_i\) 的像，记 \(\widetilde{u}_i \colon E_\pi \to F_i\) 为由 \(u_i\) 通过限制到子空间得到的 G-同构。

首先对 I 的有限子集 J 的基数作归纳，证明子空间 \((\mathrm{F}_{i})_{i\in\mathrm{J}}\) 是直和。

J 为空的情形是显然的。设 J 非空，并且所需性质对 I 中基数至多为 \(\text{Card(J)} - 1\) 的子集成立。

取 J 中的一个固定元素 \(j\)。归纳假设说明子空间 \(\mathrm{F}_{i}\)（\(i\in\mathrm{J}-\{j\}\)）构成直和。令 \(\mathrm{F}^{\prime}\) 为它们的和；现在只需证明 \(\mathrm{F}_{j}\cap\mathrm{F}^{\prime}=\{0\}\)。

假设 \(F_{j}\) 与 \(F'\) 的交不为 0；该交是 \(F_{j}\) 的一个子表示，而由于后者不可约，故有 \(F_{j} \cap F' = F_{j}\)，从而 \(F_{j} \subset F'\)。

对于 \(i \in \mathrm{J} - \{j\}\)，记 \(\operatorname{pr}_i: \mathrm{F}' \to \mathrm{F}_i\) 为投影，记 \(\iota_i: \mathrm{F}_i \to \mathrm{F}'\) 为包含映射。我们有规范同构

\[
\operatorname{Hom} _ {\mathrm{G}} \left(\mathrm{F} _ {j}, \mathrm{F} ^ {\prime}\right)\rightarrow \bigoplus_ {i \in \mathrm{J} - \{j \}} \operatorname{Hom} _ {\mathrm{G}} \left(\mathrm{F} _ {j}, \mathrm{F} _ {i}\right)
\]

（公式（1）, p. 377）使得 \(u \colon \mathrm{F}_j \to \mathrm{F}'\) 的像为族 \((\mathrm{pr}_i \circ u)_{i \in \mathrm{J} - \{j\}}\)（A, II, p. 13, cor. 1）。V, p. 387 的推论 2 表明，非零的 \(\widetilde{u}_i \circ \widetilde{u}_j^{-1}\) 是空间 \(\mathrm{Hom}_{\mathrm{G}}(\mathrm{F}_j, \mathrm{F}_i)\) 的一组基。

因此，G-态射族 \((\iota_{i} \circ \widetilde{u}_{i} \circ \widetilde{u}_{j}^{-1})_{i \in \mathrm{J} - \{j\}}\) 是 \(\mathrm{Hom}_{\mathrm{G}}(\mathrm{F}_{j}, \mathrm{F}')\) 的一组基。

记 \(\iota\) 为 \(F_{j}\) 到 \(F'\) 的包含映射；它是 \(\mathrm{Hom}_{\mathrm{G}}(\mathrm{F}_{j}, \mathrm{F}')\) 的一个元素，因而是各 G-态射 \(\iota_{i} \circ \widetilde{u}_{i} \circ \widetilde{u}_{j}^{-1}\)（\(i \in \mathrm{J} - \{j\}\)）的线性组合。于是 \(u_{j}\) 是各映射 \(u_{i}\)（\(i \neq j\)）的线性组合，这与族 \((u_{i})_{i \in \mathrm{I}}\) 的线性无关性矛盾。因此断言由归纳得到证明。

现在证明 v 是单射。设 w 是 \(\mathrm{Hom}_{\mathrm{G}}(\pi,\varrho)\otimes\mathrm{E}_{\pi}\) 的一个元素。存在唯一的有限支集族 \((x_{i})_{i\in\mathrm{I}}\)，其元素属于 \(E_{\pi}\)，使得

\[
w = \sum_ {i \in \mathrm{I}} u _ {i} \otimes x _ {i}
\]

（A, II, p. 62, cor. 1），于是有

\[
v (w) = \sum_ {i \in I} u _ {i} (x _ {i}).
\]

由前述内容，条件 \(v(w) = 0\) 因而蕴含 \(u_{i}(x_{i}) = 0\) 对所有 \(i \in I\) 成立，因而对所有 i 有 \(x_{i} = 0\)，这是因为 \(u_{i}\) 是单射，从而 w = 0。

命题 10. --- 设 G 是拓扑群。设 \(\pi\) 是 G 在复希尔伯特空间 E \(_{\pi}\) 中的不可约酉表示，\(\varrho\) 是 G 在复希尔伯特空间 E 中的酉表示。

存在 E 的闭不变子空间族 \((\mathrm{E}_{i})_{i\in\mathrm{I}}\)，使得 \(\varrho\) 在 \(E_{i}\) 中的子表示与 \(\pi\) 同构，对所有 \(i\in I\) 成立，并且 \(\mathrm{M}_{\pi}(\varrho)\) 是子空间 \(E_{i}\) 的希尔伯特直和。此外，满足这些性质的族 \((\mathrm{E}_{i})_{i\in\mathrm{I}}\) 的指标集 I 的基数与所选族无关。

证明满足所述条件的族存在。令 \(\mathcal{O}\) 为 \(\mathrm{Hom}_{\mathrm{G}}(\pi, \varrho)\setminus\{0\}\) 的子集 C 所成的集合，使得子空间 \(u(\mathrm{E}_{\pi})\)（\(u \in \mathrm{C}\)）两两正交。集合 \(\mathcal{O}\) 按包含关系排序。它非空，因为空集属于其中；它具有有限特征（E, III, p. 34, 定义 2），因为 C 属于 \(\mathcal{O}\) 当且仅当 C 的每个二元子集属于 \(\mathcal{O}\)。由 E, III, p. 35, 定理 1，\(\mathcal{O}\) 存在极大元 C。令 F 为由 \(u(\mathrm{E}_{\pi})\)（\(u \in \mathrm{C}\)）所成子空间生成的子空间的闭包；它是 \(u(\mathrm{E}_{\pi})\)（\(u \in C\)）所成空间的希尔伯特直和。我们将证明 \(\mathrm{F} = \mathrm{M}_{\pi}(\varrho)\)，从而证明族 \((u(\mathrm{E}_{\pi}))_{u \in \mathrm{C}}\) 具有所需性质。

按定义，\(u(\mathrm{E}_{\pi}) \subset \mathrm{M}_{\pi}(\varrho)\) 对所有 \(u \in \mathrm{C}\) 成立，故 \(\mathrm{F}\) 包含于 \(\mathrm{M}_{\pi}(\varrho)\)。为证明反向包含，只需证明：若 \(v\) 是从 \(\pi\) 到 \(\varrho\) 的 G-态射，则其像包含于 \(\mathrm{F}\)。令 \(p\) 为 \(\mathrm{E}\) 上像为 \(\mathrm{F}^{\circ}\) 的正交投影；由于 \(\mathrm{F}\) 是 \(\mathrm{E}\) 的子表示，它是 G-态射（V, p. 383 的命题 4）。G-态射 \(p \circ v\) 的像与 \(\mathrm{F}\) 正交；因此该像为零（否则 \(\mathrm{C} \cup \{p \circ v\} \in \mathcal{O}\)，这与 \(\mathrm{C}\) 的极大性矛盾），从而 \(v\) 的像包含于 \(\mathrm{F}\)。

现在设 \((\mathrm{E}_{i})_{i\in\mathrm{I}}\) 和 \((\mathrm{F}_{j})_{j\in\mathrm{J}}\) 是 E 的两两正交闭不变子空间族，使得 G 在 \(E_{i}\)（分别在 \(F_{j}\)）中的子表示与 \(\pi\) 同构，对所有 \(i\in I\)（分别对每个 \(j\in J\)）成立，并且 \(\mathrm{M}_{\pi}(\varrho)\) 既是族 \((\mathrm{E}_{i})_{i\in\mathrm{I}}\) 的希尔伯特直和，也是族 \((\mathrm{F}_{j})_{j\in\mathrm{J}}\) 的希尔伯特直和。

若 I 有限，则

\[
\dim \operatorname{Hom} _ {\mathrm{G}} (\pi , \mathrm{M} _ {\pi} (\varrho)) = \dim \operatorname{Hom} _ {\mathrm{G}} \Bigl (\mathrm{E} _ {\pi}, \bigoplus_ {i \in \mathrm{I}} \mathrm{E} _ {i} \Bigr) = \operatorname{Card} (\mathrm{I})
\]

（V, p. 377 的公式（1）和 V, p. 387 的推论 2）。于是对于 J 的每个有限子集 L，有

\[
\begin{array}{c} \operatorname{Card} (\mathrm{L}) = \dim \operatorname{Hom} _ {\mathrm{G}} \Bigl (\mathrm{E} _ {\pi}, \bigoplus_ {j \in \mathrm{L}} \mathrm{F} _ {j} \Bigr) \\ \leqslant \dim \operatorname{Hom} _ {\mathrm{G}} (\pi , \operatorname{M} _ {\pi} (\varrho)) = \operatorname{Card} (\mathrm{I}) \end{array}
\]

（同上）。这证明了 J 也是有限的，且 \(\text{Card(I)} = \text{Card(J)}\)，正如所要求的。同样，若 J 有限，则 I 有限且具有相同的基数。

现在假设 I 和 J 都是无限的。对于 \(j \in J\)，记 \(p_{j}\) 为 E 上像为 \(F_{j}\) 的正交投影。对于 \(i \in I\)，取 \(x_{i}\) 为 \(E_{i}\) 中的非零元素；它是 \(E_{i}\) 的循环向量。注意，由于 \(p_{j}\) 通过限制到子空间得到的是从 \(E_{i}\) 到 \(F_{j}\) 的 G-态射，空间 \(p_{j}(E_{i})\) 为零，当且仅当 \(p_{j}(x_{i}) = 0\)（事实上，空间 \(E_{i} \cap \text{Ker}(p_{j})\) 要么为零，要么等于 \(E_{i}\)）。

对于每个 \(j \in J\)，存在 \(f(j) \in I\)，使得 \(p_{j}(\mathrm{E}_{f(j)})\) 不为 0（否则，正交投影 \(p_{j}\) 在 \(\mathrm{M}_{\pi}(\varrho)\) 上为零）。于是定义了从 J 到 I 的映射 f。对于每个 \(i \in I\)，集合 \(f^{-1}(i)\) 是可数的。事实上，该集合包含于由 \(j \in J\) 中满足 \(p_{j}(\mathrm{E}_{i})\) 非零者组成的集合中，也就是满足 \(p_{j}(x_{i}) \neq 0\) 的那些 j。现在由（EVT, V, p. 20, cor. 2），有

\[
\sum_ {j \in \mathrm{J}} \| p _ {j} (x _ {i}) \| ^ {2} = \| x _ {i} \| ^ {2},
\]

因此，\(j \in J\) 中满足 \(p_{j}(x_{i}) \neq 0\) 的元素所成集合是可数的。由 E, III, p. 50, 命题 4，得到 \(\text{Card}(\text{J}) = \text{Card}(f(\text{J})) \leqslant \text{Card}(\text{I})\)。交换 I 和 J 的角色，得到 \(\text{Card}(\text{I}) = \text{Card}(\text{J})\)。

推论. --- 设 G 是拓扑群，\(\varrho\) 是 G 在复希尔伯特空间 E 中的酉表示。G 的 \(\pi\)-同型分量（其中 \(\pi \in \widehat{\mathbf{G}}\)）的希尔伯特直和是 \(\varrho\) 的最大半单子表示。

事实上，\(\pi\)-同型分量对每个 \(\pi \in \widehat{\mathbf{G}}\) 两两正交（命题 8，\(a\)），且每个都是半单的（命题 10，\(a\)），所以空间 \(\mathrm{M}_{\pi}(\varrho)\) 对 \(\pi \in \widehat{\mathbf{G}}\) 的希尔伯特直和 F 定义了 \(\varrho\) 的一个半单子表示。此外，\(\varrho\) 的每个不可约子表示都是 \(\varrho\) 的某个同型分量的子表示，故推论得证。

定义 10. --- 设 G 是拓扑群。设 \(\varrho\) 是 G 在复希尔伯特空间 E 中的酉表示，\(\pi\) 是 G 在复希尔伯特空间中的不可约酉表示。

称 \(\pi\) 在 \(\varrho\) 中的重数为指标集的基数，所针对的是 E 的在 G 下稳定的闭子空间族 \((\mathrm{E}_i)_{i\in \mathrm{I}}\)，其中 \(\varrho\) 在 \(\mathrm{E}_i\) 中诱导的子表示与 \(\pi\) 同构，对每个 \(i\in I\) 成立，并且 \(\mathrm{M}_{\pi}(\varrho)\) 是子空间 \(\mathrm{E}_i\) 的希尔伯特直和。

如果 \(\pi\) 在 \(\varrho\) 中的重数有限，则由 V, p. 377 的公式（1）和 V, p. 387 的推论 2，它等于空间 \(\mathrm{Hom}_{\mathrm{G}}(\pi, \varrho)\) 的维数（分别等于 \(\mathrm{Hom}_{\mathrm{G}}(\varrho, \pi)\) 的维数）。一般而言这并不总是成立。

注. --- 一个酉表示 \(\varrho\) 可能非零，但相对于 G 的所有不可约表示的 \(\varrho\) 的所有同型分量都为零（参见 V, p. 426, 注）。

